\section{Conclusion}\label{s:conclusion}

In this thesis, we evaluated the performance and ranking of six ABR algorithms in Mahimahi and CellReplay. In both emulators, we replayed the same 5G traces to isolate emulator differences. We computed the average results for several metrics to reflect the overall performance of the ABR algorithms in both emulators. In our findings, we observed that the SSIM (video quality) and BufRatio (playback stability) rankings remain the same, but the values improve slightly when CellReplay is used. In contrast, we noticed significant changes in the RTT and delivery rate values. CellReplay provides around 60\% increase in delivery rate, and much more stable RTT values. Buffer time rankings slightly shift when moving from Mahimahi to CellReplay, but the differences in values are low. We also reassessed the trace variability by replaying 4G and 5G traces in Mahimahi, and saw significant improvements in all aspects, reflecting how ABR algorithms behave under different cellular generations. Our experiments have some limitations. For example, we evaluate a single 10-minute video and use only cellular traces. Future work could extend this by testing longer videos with higher quality and replaying other network types such as WiFi or satellite. We are using pre-trained learning-based algorithms that are not tuned for modern 5G cellular networks. For future experiments, training these algorithms according to the test environment can help improve their performance. Lastly, CellReplay provides a realistic emulation environment, but real-world deployments still involve additional complexities beyond our scope.

%%% Local Variables:
%%% mode: latex
%%% TeX-master: "../thesis"
%%% End:
